% First Week Survival Guide
% Auto-generated from megn300_master_reference.tex

\chapter*{First Week Survival Guide}
\addcontentsline{toc}{chapter}{First Week Survival Guide}

If this is your first day in MEGN~300, start here:

\begin{enumerate}[leftmargin=2em]
\item \textbf{Read Chapter~1} (Fundamentals of Measurement Theory). Focus on the five-block system and sensor terminology.
\item \textbf{Read Chapter~\ref{ch:debugging}} (Systematic Debugging). Read it \emph{before} you build your first circuit, not after it fails.
\item \textbf{Read the LabVIEW appendix}. Familiarize yourself with DAQmx, essential VIs, and the ``Common Mistakes'' checklist.
\item \textbf{Review the Laboratory Safety chapter} (page~\pageref{ch:safety}). You cannot enter the lab without understanding electrical safety.
\item \textbf{Practice}: open LabVIEW, drop a DAQmx Read VI, wire it to a Waveform Chart in a While Loop with a 100\,ms Wait, and run it with the NI~USB-6001. If you can read one voltage channel, you are ready for Lab~1.
\end{enumerate}

\begin{warningbox}[Do Not Read This Document Cover to Cover]
This is a 200+ page reference. Use the ``What Should I Be Reading Right Now?'' triage guide to find exactly the section you need for each module.
\end{warningbox}

\begin{tipbox}[Also Read Before Your First Lab]
The ``Why Instrumentation Projects Fail'' chapter (page~\pageref{ch:projects_fail}) consolidates the ten most common failure modes from six years of MEGN~300 projects. Ten minutes of reading it now saves hours of debugging later.
\end{tipbox}

\newpage
